\section{Announcements}
\begin{itemize}
    \item Pset \#11 is due on Wednesday November 18.
    \item Pset \#12: Sampling will be released on Wednesday November 18 and is due on Wednesday December 2.
    \item Midterm exam \#3 will be on Wednesday December 2.
    \item Nothing is due on November 25 (Thanksgiving week).
    \item \textbf{Wikipedia assignment:} ``Start drafting your contributions'' is due
    \textbf{\mbox{Sunday November 22}}. Complete the ``Drafting as a group'' and
    ``Plagiarism'' training modules by the same date; see the
    \href{https://dashboard.wikiedu.org/courses/UC_Merced/Signal_Processing_and_Linear_Systems_(Fall_2026)/timeline}{Wiki Education timeline}.
\end{itemize}

\section{Filtering to remove distortions}

The original continuous-time cosine only had two impulses 
at $\omega = \pm\omega_0$. After sampling, $X_s(\omega)$ contains infinitely
many impulses, at
\[
\omega = \pm\omega_0, \quad
\pm\omega_0 \pm \omega_s, \quad
\pm\omega_0 \pm 2\omega_s, \quad \dots
\]

If $\omega_s$ is large enough, these copies are separated in frequency and do not overlap.
But if $\omega_s$ is too small, the shifted impulses can collide or cross into each other.
In other words, impulses from different ``copies'' of $X(\omega)$ may end up in the same
frequency location. This leads to distortion in the reconstructed signal called aliasing.

\textbf{What is aliasing?} When the sampling frequency is too low,
different frequency components in the original signal produce overlapping or indistinguishable
contributions in the sampled spectrum. In the reconstructed signal, these contributions
appear as \emph{fake} or \emph{misplaced} frequencies. This phenomenon is called \textbf{aliasing}:
different original frequencies become indistinguishable after sampling.

For a single cosine, aliasing means that a high-frequency cosine can ``look like'' a lower-frequency cosine once sampled, because their impulses line up after shifts
by $\omega_s$.


As we said above, to reconstruct the original continuous-time signal $x(t)$ from its samples, we may design a low-pass filter to remove the extra copies in $X_s(\omega)$. This will work only if the copies do not overlap with the original spectrum. An ideal low-pass filter has a cutoff frequency $\omega_c$ that passes only the frequency band containing the original $X(\omega)$ (from $-\omega_c$ to $+\omega_c$) and removes all higher frequencies will work well.

For a single cosine at $\omega_0$, we can choose the cutoff frequency $\omega_c$ such that
\[
\omega_0 < \omega_c < \omega_s - \omega_0.
\]
This is possible \emph{only if} the original impulse at $+\omega_0$ is strictly
separated from the nearest shifted impulse at $\omega_s - \omega_0$. That gives the
inequality
\[
\omega_0 < \omega_s - \omega_0
\quad \Longleftrightarrow \quad
\omega_s > 2\omega_0.
\]

Under this condition, an ideal low-pass filter with passband
$|\omega| < \omega_c$ will keep the original impulses at $\pm\omega_0$ and completely remove all impulses from the shifted copies.
In the time domain, this low-pass filter reconstructs the original continuous-time cosine
$x(t) = \cos(\omega_0 t)$ from the samples (our overall goal!). The condition that we just obtained is called the \textbf{Nyquist sampling condition}.

We have shown that to avoid overlap between the original impulses at $\pm\omega_0$
and the nearest replicated impulses (at $\pm\omega_0 \pm \omega_s$), we must have
\[
\omega_s > 2\omega_0.
\]

We can obtain an equivalent relation in Hz. Recall $\omega_s = \dfrac{2\pi}{T_s}$ and $\omega_0 = 2\pi f_0$ where $f_0$ is the
ordinary frequency in Hz. In terms of the sampling frequency $F_s = 1/T_s$:
\[
\omega_s > 2\omega_0
\quad\Longleftrightarrow\quad
2\pi F_s > 2(2\pi f_0)
\quad\Longleftrightarrow\quad
F_s > 2f_0.
\]

So, the Nyquist rate for a single cosine:
\[
F_s > 2f_0, \qquad f_0 = \frac{\omega_0}{2\pi}.
\]
If we sample \emph{above} this Nyquist rate, we can in principle use an ideal
low-pass filter to remove the extra copies in $X_s(\omega)$ and reconstruct the
original cosine exactly from its samples. This is the key result of this week's theoretical topics.

\section{Nyquist Sampling Theorem}
We end with the formal statement of the Nyquist Sampling Theorem for general bandlimited signals (not just a single cosine).
\begin{theorem}[Nyquist Sampling Theorem]
If a continuous-time signal $x(t)$ contains no frequencies higher than
$B$ Hz, that is, if its Fourier transform $X(\omega)$ is zero for
$|\omega| > 2\pi B$, then $x(t)$ is completely determined by its samples
taken at a sampling frequency $F_s$ greater than $2B$ Hz. In this case,
$x(t)$ can be reconstructed from its samples using an ideal low-pass filter. 
\end{theorem}

\section{Recommended Reading}
From the week 12 folder, try the sampling app shown in class by running \texttt{streamlit run sampling\_app/sampling\_app.py}. This will let you play with different sampling frequencies and low-pass cut off frequencies to see its effects. Also, run \href{https://github.com/ee-ucmerced/ee102-signals-systems/blob/main/lecture_notes/week12_sampling/VM_aliasing.ipynb}{\texttt{VM\_aliasing.ipynb}} notebook that is available on course GitHub to apply these concepts to the guitar tune example. Try to achieve the following using the code:
\begin{itemize}
    \item Visualize aliasing with $x(t) = \sin(2\pi f_0 t)$ for different values of $f_0$ and different sampling frequencies $F_s$.
    \item Observe how increasing the number of samples improves the reconstruction of the original signal for the sine wave.
    \item Apply the sampling theorem to adequately sample and reconstruct your recorded voice from Problem Set \#10.
\end{itemize}
